\section*{Chapitre \ref{ch:s2:convertible-arbitrage} --- Arbitrage de convertibles}

\begin{solution}{exo:s2:convertible-arbitrage:1}
$100/2.5 = 40$ par action~; à 60, la valeur de conversion vaut $2.5 \times 60 = 150$.
\end{solution}

\begin{solution}{exo:s2:convertible-arbitrage:2}
Ses obligations représentent trois fois ses fonds propres, soit $3 \times 2.7 = 8.1~\%$ des fonds propres, ce qui est proche des $-7.2~\%$ publiés par les fonds pour la période de janvier à mai 2005.
\end{solution}

\begin{solution}{exo:s2:convertible-arbitrage:3}
$1.27 \times 15/81.2 = 0.23$~: une baisse de 1~\% de l'action retire 0.23~\% à l'obligation.
\end{solution}

\begin{solution}{exo:s2:convertible-arbitrage:4}
L'obligation simple est actualisée au spread de crédit correspondant à l'intensité de défaut de l'émission, pour un cours de l'action de 40. À 15, le lien action-crédit a relevé
l'intensité de défaut~; l'obligation vaut donc moins qu'une obligation au spread initial, et l'option de conversion est trop hors de la monnaie pour compenser cette baisse.
\end{solution}

\begin{solution}{exo:s2:convertible-arbitrage:5}
Convexité~: lorsque l'action a baissé, le delta de l'obligation a diminué, de sorte qu'elle a perdu moins par unité de baisse que la couverture vendeuse, dimensionnée au delta plus élevé, n'a gagné. Le carnet
est long en gamma, et les fortes variations du krach ont rapporté 3.7 points nets.
\end{solution}

\begin{solution}{exo:s2:convertible-arbitrage:6}
La cheapness est l'écart entre le prix de marché et la valeur du modèle, qui dépend de ceux qui doivent vendre et de ceux qui peuvent acheter. L'action, le credit default swap (CDS) et le
marché des options ne la portent pas, donc aucune couverture négociée ne la compense~; seule la capacité à tenir la position jusqu'à sa résorption le fait.
\end{solution}

\begin{solution}{exo:s2:convertible-arbitrage:7}
Sur cinq ans, la couverture en crédit coûte 21.8~\% des fonds propres en primes et abandonne le rebond du crédit~; elle ramène le carry de 9.8~\% à 6.0~\% par an, et le rendement sur cinq ans de 51.9~\% à 30.1~\%.
Elle absorbe les 6.9 points de perte de crédit de l'épisode, et réduit de moitié sa perte, de $-11.6~\%$ à
$-5.8~\%$, et laisse intacts les 12.6 points de cheapness.
\end{solution}

\begin{solution}{exo:s2:convertible-arbitrage:8}
Les couvertures en delta et en crédit laissent subsister le risque de cheapness, celui de l'emprunt et celui du financement~: dans l'épisode du chapitre, la seule cheapness a coûté 12.6~\% des fonds propres à un levier de trois pour un,
de sorte qu'à dix pour un elle coûterait environ 42~\%, avant tout appel de marge. En 2005 comme en 2008, le risque qui comptait était précisément celui que les couvertures laissaient en place.
\end{solution}

\begin{solution}{pb:s2:convertible-arbitrage:1}
\begin{enumerate}
\item Une obligation simple, une option sur les actions, moins le call de l'émetteur, le tout exposé au défaut.
\item 117.7 contre une valeur de conversion de 100, un delta de 1.80 actions.
\item Une convertible très en dessous de son prix de conversion, qui se traite sur le crédit~; à 15, elle vaut 81.2, en dessous de l'obligation simple à 87.34, avec un delta de 1.27
et une élasticité de 0.23.
\item Le delta de la couverture inclut l'effet crédit du cours de l'action~; une couverture qui l'ignore est trop faible lorsque l'action baisse.
\item Selon la définition du chapitre~: une position vendeuse sur l'action égale au delta, plus une protection de crédit dimensionnée à la sensibilité au crédit.
\item Une stratégie d'achat et de couverture explique une bonne part des rendements des fonds~; les événements extrêmes et l'offre de convertibles comptent aussi.
\item Des rachats de plus de 20~\% en un trimestre, des ventes forcées, une décote médiane de 2.7~\% en mai, des rendements de fonds de $-7.2~\%$.
\item L'effet de levier offert par les prime brokers s'est brusquement réduit lorsque ceux-ci ont redouté la défaillance.
\item Trente obligations à l'émission, couvertes en delta chaque jour, à un levier de trois pour un~; P\&L en gamma, en crédit, en cheapness, en coupons et en emprunt.
\item 9.8~\% par an couverts en delta~; 6.0~\% couverts en crédit.
\item $-11.6~\%$~: cheapness $-12.6$, crédit $-6.9$, gamma $+3.7$, coupon $+5.3$, emprunt $-1.1$.
\item Les obligations ont retrouvé leur cheapness et l'intensité de défaut a reflué~: $+22.0~\%$ en six mois.
\item \textbf{Résultat nommé.} Le carnet couvert en delta rapporte 9.8~\% par an sur les fonds propres et perd 11.6~\% dans l'épisode trimestriel simulé, dont 6.9 points dus au
crédit et 12.6 à la cheapness~; couvrir le crédit réduit de moitié la perte à 5.8~\% au prix du carry (6.0~\% par an).
\item Elle coûte 21.8~\% des fonds propres sur cinq ans et absorbe les 6.9 points de perte de crédit.
\item La cheapness est multipliée par le levier et peut forcer des ventes aux pires prix.
\item Ils ont fait face à des rachats (2005) ou ont perdu leur effet de levier (2008)~: le capital, et non les prix, a tranché.
\item Par la perte dans un épisode de ventes forcées au levier retenu, à comparer au capital qui pourrait être perdu sans débouclage forcé.
\item Des valeurs de modèle avec des données réalistes, des épisodes de cheapness, et des coûts d'emprunt et de financement qui varient.
\item La convertible nouvellement émise.
\item Pour financer rapidement des émetteurs et conserver des obligations peu liquides pendant les périodes où d'autres doivent vendre.
\end{enumerate}
\end{solution}

\begin{solution}{iq:s2:convertible-arbitrage:1}
Acheter la convertible, vendre à découvert le delta en actions, éventuellement acheter une protection de crédit~; percevoir le coupon, la rémunération de la vente à découvert diminuée de l'emprunt, le gamma
issu du rééquilibrage et la convergence d'une obligation bon marché vers sa valeur~; les risques sont la cheapness, le crédit, l'emprunt et le financement.
\end{solution}

\begin{solution}{iq:s2:convertible-arbitrage:2}
Lorsque l'action baisse, l'intensité de défaut augmente, si bien que l'obligation chute davantage qu'un modèle d'option à crédit constant ne le prévoit~: le delta est plus grand aux faibles cours de l'action,
et, aux cours très bas, l'essentiel du delta relève du crédit.
\end{solution}

\begin{solution}{iq:s2:convertible-arbitrage:3}
La cheapness, le rappel de l'emprunt, le retrait de financement, le gamma dans les sauts de cours, les calls de l'émetteur et les offres publiques, l'erreur de modèle dans l'entrée de crédit~; tester avec un scénario de cheapness de type 2005,
une intensité de défaut doublée et un rappel de l'emprunt simultanés, au levier du carnet.
\end{solution}

\begin{solution}{iq:s2:convertible-arbitrage:4}
Trouver de l'emprunt ailleurs, couvrir avec des options ou un instrument corrélé, ou réduire la position~; accepter une couverture moins précise plutôt que
d'une obligation non couverte, et consigner le coût pour le risque d'emprunt du carnet.
\end{solution}

\begin{solution}{iq:s2:convertible-arbitrage:5}
Chaque jour~: cours des obligations, des actions, du crédit et des volatilités~; valeurs de modèle et deltas issus du pricer (tabulés pour la rapidité)~; ajustements de couverture compensés entre les
émetteurs~; disponibilité de l'emprunt~; suivi de la cheapness et attribution~; contrôles des opérations sur titres, des rachats et des calls et des conversions.
\end{solution}

\begin{solution}{iq:s2:convertible-arbitrage:6}
Comme pour toute option, la valeur de la position couverte varie de $\Theta\,\mathrm{d}t + \tfrac12\Gamma(\mathrm{d}S)^2$, auxquels s'ajoutent comme carry le coupon de l'obligation et la
rémunération de la vente à découvert diminuée de l'emprunt. Au défaut, l'action fait un saut vers une valeur proche de zéro~: la vente à découvert de l'action gagne sa valeur, et l'obligation tombe à son taux de recouvrement~; avec un modèle action-crédit,
la couverture en delta a été dimensionnée pour une partie de ce saut, et la protection de crédit couvre le reste.
\end{solution}
